Los resultados experimentales corroboran que la notación asintótica $\mathcal{O}$ describe la tendencia de crecimiento hacia el infinito, pero no anticipa el desempeño efectivo sobre arquitecturas modernas en instancias acotadas. En ordenamiento, \texttt{std::sort} y QuickSort dominan debido a su mínima sobrecarga auxiliar y óptima localidad espacial. En multiplicación de matrices, la formulación directa de Strassen demostró ser inviable frente a Naive (tardando 152 segundos versus 0.15 segundos para $N=1024$), evidenciando que las constantes ocultas, el tráfico hacia la memoria principal y la falta de un corte híbrido hacia métodos iterativos pueden degradar severamente algoritmos de menor complejidad asintótica.